% from ega4-2-fr.tex
\oldpage[IV-2]{19}
(II, 7.4.2), pour $X$ le préschéma somme des schémas
$\operatorname{Spec}(\mathcal O_y)$, où $y$ parcourt $Y$ (I, 3.1), pour $f$ le
morphisme canonique, qui est fidèlement plat (I, 2.4.2)~; si $\eta$ désigne le
point générique de $Y$, $Z=\{\eta\}$ n'est pas ouvert dans $Y$ (II, 7.4.3),
mais $f^{-1}(Z)=\operatorname{Spec}(\mathcal O_\eta)$ puisque
$\mathcal O_\eta$ est un corps, et par suite $f^{-1}(Z)$ est ouvert dans $X$.
\end{remark}

\subsection{Morphismes universellement ouverts et morphismes plats}

\begin{env}[2.4.1]
\label{IV.2.4.1-fr}
Nous avons déjà défini (II, 5.4.9) la notion de morphisme
\emph{universellement fermé}~; de la même manière, on pose les définitions
suivantes~:
\end{env}

\begin{definition}[2.4.2]
\label{IV.2.4.2-fr}
On dit qu'un morphisme de préschémas $f:X\to Y$ est \emph{universellement
ouvert} (resp. \emph{universellement bicontinu}, resp. un
\emph{homéomorphisme universel}) si, pour tout morphisme $g:Y'\to Y$, le
morphisme $f_{(Y')}:X\times_Y Y'\to Y'$ est ouvert (resp. un homéomorphisme
sur son image, resp. un homéomorphisme sur $Y'$).
\end{definition}

Nous verrons plus loin (14.3.2) que lorsque $Y$ est localement noethérien, la
définition de morphisme universellement ouvert donnée ici est équivalente à la
définition (III, 4.3.9) pour les morphismes de type fini~; le lecteur pourra
vérifier que nous n'utilisons pas cette dernière définition avant le \S~14.

\begin{proposition}[2.4.3]
\EGAReaderTarget{ega4.c04.prop.2-4-3}{IV.2.4.3-fr}
\begin{enumerate}
\item[(i)] Une immersion (resp. une immersion ouverte, resp. une immersion
fermée) est universellement bicontinue (resp. universellement ouverte, resp.
universellement fermée).
\item[(ii)] Le composé de deux morphismes universellement ouverts (resp.
universellement fermés, resp. universellement bicontinus, resp. deux
homéomorphismes universels) l'est aussi.
\item[(iii)] Si $f:X\to Y$ est un $S$-morphisme universellement ouvert (resp.
universellement fermé, resp. universellement bicontinu, resp. un
homéomorphisme universel), il en est de même de
$f_{(S')}:X_{(S')}\to Y_{(S')}$ pour tout changement de base $S'\to S$.
\item[(iv)] Si $f:X\to X'$, $g:Y\to Y'$ sont deux $S$-morphismes
universellement ouverts (resp. universellement fermés, resp. universellement
bicontinus, resp. deux homéomorphismes universels), il en est de même de
\[
f\times_S g:X\times_S Y\longrightarrow X'\times_S Y'.
\]
\item[(v)] Soient $f:X\to Y$, $g:Y\to Z$ deux morphismes tels que $f$ soit
surjectif~; si $g\circ f$ est universellement ouvert (resp. universellement
fermé, resp. universellement bicontinu, resp. un homéomorphisme universel), il
en est de même de $g$.
\item[(vi)] Pour que $f$ soit un morphisme universellement ouvert (resp.
universellement fermé, resp. universellement bicontinu, resp. un
homéomorphisme universel), il faut et il suffit que $f_{\mathrm{red}}$ le soit.
\item[(vii)] Soit $(U_\alpha)$ un recouvrement ouvert de $Y$. Pour que
$f:X\to Y$ soit universellement ouvert (resp. universellement fermé, resp.
universellement bicontinu, resp. un homéomorphisme universel), il faut et il
suffit que pour tout $\alpha$, sa restriction
$f^{-1}(U_\alpha)\to U_\alpha$ le soit.
\end{enumerate}
\end{proposition}

L'assertion (i) résulte de (I, 4.3.2). L'assertion (ii) résulte aussitôt des
définitions, et il en est de même de (iii) en se ramenant au cas où $Y=S$,
$Y'=S'$, ce que l'on peut faire grâce à (I, 3.3.11)~; on sait que (iv)
découle de (ii) et (iii) (I, 3.5.1). Pour prouver (v), notons que pour tout
morphisme $Z'\to Z$, $f_{(Z')}:X_{(Z')}\to Y_{(Z')}$ est surjectif
(I, 3.5.2)~; on peut donc se borner à prouver que si $g\circ f$ est ouvert
(resp. fermé, resp. un
